
\subsubsection*{5.2.1}
\phantomsection
\label{sga3:VIIB:5.2.1}

Let \(A\) be the affine algebra of \(G\), let \(I_A\) be its
augmentation ideal, and put
\[
  I=B\cap I_A.
\]
By \SGARefLink{sga3:VIIB:2.4}{2.4},
\[
  H=\operatorname{Spf}\bigl(A/\overline{AI}\bigr)
\]
and
\[
  B
  =
  \left\{
    x\in A\ \middle|\
    \Delta(x)-x\widehat\otimes1
    \in A\widehat\otimes\overline{AI}
  \right\}.
\]
\phantomsection
\label{sga3:VIIB:5.2.1:tensor-order-source-note}
\noindent\textit{Source note.}\enspace
The French PDF prints
\(\Delta(x)-1\widehat\otimes x\) here. The right-quotient
convention fixed in \SGARefLink{sga3:VIIB:5.0}{5.0} and \SGARefLink{sga3:VIIB:5.1:proposition}{Proposition~5.1}, together with the
defining condition
\((\operatorname{id}_A\widehat\otimes\pi)\Delta(x)
=x\widehat\otimes1\), requires
\(\Delta(x)-x\widehat\otimes1\), which is used above.

Put
\[
  \omega=I/\overline{I^2}.
\]
For each \(r\), denote by \(I_r\) the closed subideal of \(I\)
consisting of the elements \(x\) such that
\[
  x^{p^r}=0,
\]
and denote by \(\omega_r\) the canonical image of \(I_r\) in
\(\omega\). We shall prove the following proposition.

\medskip
\noindent\textbf{Proposition.}
\phantomsection
\label{sga3:VIIB:5.2.1:section-proposition}
{\itshape If there exists a continuous section
\[
  \sigma:\omega\longrightarrow I
\]
of the projection \(I\to I/\overline{I^2}\) such that
\[
  \sigma(\omega_r)\subset I_r
\]
for every \(r\), then \(B\) is isomorphic to
\[
  k[[\omega]]
  \big/
  \bigl((x^{p^r})\bigr)_{x\in\omega_r}.
\]}

Indeed, such a section extends to a continuous morphism
\[
  k[[\omega]]\longrightarrow B,
\]
which factors through
\[
  B'
  =
  k[[\omega]]
  \big/
  \bigl((x^{p^r})\bigr)_{x\in\omega_r}.
\]
In \SGARefLink{sga3:VIIB:5.2.2}{5.2.2}--\SGARefLink{sga3:VIIB:5.2.5}{5.2.5} we prove that the resulting morphism
\[
  \varphi:
  B'
  =
  k[[\omega]]
  \big/
  \bigl((x^{p^r})\bigr)_{x\in\omega_r}
  \longrightarrow B
\]
is an isomorphism.
